\chapter{Grigori Perelman}\label{ch2006grigoriperelmandeclined:grigori-perelman-fields-medal-2006-declined}

\begin{center}\rule{0.5\linewidth}{0.5pt}\end{center}

\textbf{Prepared:} July 2026 \textbf{Subject:} Grigori Yakovlevich Perelman (b. June 13, 1966, Leningrad, USSR) \textbf{Position:} Former researcher, Steklov Institute of Mathematics, St.~Petersburg \textbf{Doctoral Advisor:} Yuri Burago (Steklov Institute) \textbf{Fields Medal Citation:} \emph{``For his contributions to geometry and his substantial insights into the analytical and geometric structure of the Ricci flow.''}

\begin{center}\rule{0.5\linewidth}{0.5pt}\end{center}

\begin{prereqs} Riemannian geometry (curvature tensors, comparison geometry), PDE theory (parabolic equations, Hamilton--Jacobi equations), basic 3-manifold topology (fundamental groups, surgery theory). \end{prereqs}

\begin{objectives} Explain the Ricci flow with surgery and its role in proving the Poincar\'e conjecture; construct the Perelman entropy functional and prove its monotonicity; describe the no-local-collapsing theorem and the classification of singularity models; understand finite extinction time. \end{objectives}

\section{Problem History}

The Poincar\'e conjecture, formulated by Henri Poincar\'e in 1904, asks whether every simply connected closed 3-manifold is homeomorphic to the 3-sphere $S^3$. Despite its simple statement, the problem resisted every attempt at proof for nearly a century. Stephen Smale proved the conjecture for dimensions $n \ge 5$ (Fields Medal 1966), and Michael Freedman settled the 4-dimensional case (Fields Medal 1982), but dimension 3---the original and hardest case---remained open.

In the late 1970s, William Thurston proposed his geometrization conjecture: every closed orientable 3-manifold can be decomposed into pieces, each carrying one of eight model geometries (spherical, Euclidean, hyperbolic, $S^2 \times \mathbb{R}$, $\mathbb{H}^2 \times \mathbb{R}$, $\operatorname{SL}_2(\mathbb{R})$, Nil, Sol). This conjecture implies the Poincar\'e conjecture as a special case. Thurston proved it for Haken manifolds (Fields Medal 1982), but the general case---including non-Haken manifolds---remained open.

Richard Hamilton introduced the Ricci flow in 1982 as a tool for attacking geometrization. The Ricci flow evolves a Riemannian metric according to $\partial_t g_{ij} = -2R_{ij}$, where $R_{ij}$ is the Ricci curvature tensor. The idea is that the flow diffuses curvature, smoothing out irregularities. Hamilton proved short-time existence, convergence for positive Ricci curvature, and classified singularity types, but several crucial steps remained beyond his reach: controlling singularities, showing finite extinction time, and making the surgery construction precise. The Ricci flow is a nonlinear heat equation for the metric, and it is \emph{quasilinear parabolic}: the principal symbol is the Laplace--Beltrami operator of the evolving metric. The short-time existence is proved using the DeTurck trick, which modifies the flow by a diffeomorphism to make it strictly parabolic, and then solves the modified equation using standard PDE methods.

A central obstacle was the \emph{pinching problem}: as the Ricci flow evolves, regions of high curvature can develop where the ratio of the maximum to minimum sectional curvature becomes unbounded. Hamilton developed a powerful maximum principle to control this ratio under certain conditions, but the general case---especially near singularities---required Perelman's new entropy methods. Hamilton's pinching theorem states that if the initial metric satisfies a curvature pinching condition of the form $\operatorname{Rm} \ge \varepsilon R \cdot g$ for some $\varepsilon > 0$, then the pinching improves under the flow and the solution exists for all time and converges to a metric of constant sectional curvature. However, this condition is not satisfied in general, and Perelman's entropy methods were needed to control the curvature ratio near singularities.

Hamilton's program had four main components, each of which required Perelman's insights: (1) short-time existence and uniqueness of the Ricci flow with surgery, (2) the canonical neighborhood theorem describing the geometry of high-curvature regions, (3) the no-local-collapsing theorem preventing degenerate singularities, and (4) the finite extinction time argument. Perelman's genius was to introduce new monotone functionals---the $\mathcal{F}$-functional, the $\mathcal{W}$-functional, and the $\mathcal{L}$-length---that provided the missing tools for each component.

\section{How Was It Solved}

In three preprints posted on the arXiv in 2002--2003 \cite{perelman2002}, totaling about 130 pages, Perelman \cite{perelman2002} filled in all the missing steps of Hamilton's program.

\textbf{The entropy functional.} Perelman \cite{perelman2002} introduced the functional:
\[
\mathcal{W}(g, f, \tau) = \int_M \left[ \tau(R + |\nabla f|^2) + f - n \right] (4\pi\tau)^{-n/2} e^{-f} \, dV
\]
and proved its monotonicity: $\frac{d}{dt} \mathcal{W} \ge 0$ under the coupled Ricci flow and backward heat equation for $f$. The $\mathcal{F}$-functional $\mathcal{F}(g,f) = \int_M (R + |\nabla f|^2) e^{-f} \, dV$ is the gradient flow analogue of the Dirichlet energy. The monotonicity of $\mathcal{F}$ under the coupled Ricci flow and the backward heat equation $\partial_t f = -R - |\nabla f|^2$ was proved using the evolution formula:
\[
\frac{d}{dt} \mathcal{F}(g(t), f(t)) = \int_M \left| \operatorname{Rc} + \nabla^2 f \right|^2 e^{-f} \, dV \ge 0.
\]
This identity, sometimes called the \emph{Bochner formula for the Ricci flow}, is the engine behind Perelman's monotonicity results.

\textbf{No-local-collapsing theorem.} Using the entropy functional, Perelman proved that the Ricci flow cannot develop local collapsing: if the curvature is bounded at a point, then the volume of a small ball around that point cannot be too small. This prevents degenerate singularities and ensures that any singularity must be of a special type---a neck pinch where the manifold looks like $S^2 \times \mathbb{R}$. The theorem is sharp: the ratio $\mathcal{W}(g,f,\tau) / \mathcal{F}(g,f)$ provides a lower bound on the volume of geodesic balls, and the no-local-collapsing property is used crucially in the canonical neighborhood theorem.

\textbf{Surgery construction.} Perelman developed a precise method for performing surgery: when a singularity develops, identify the high-curvature region as a neck (diffeomorphic to $S^2 \times I$), cut along the neck, cap the boundaries with 3-balls, and restart the flow. The no-local-collapsing theorem and the classification of singularity models ensure that only finitely many surgeries occur before the manifold becomes extinct. Perelman proved that each surgery reduces the Euler characteristic by at most a bounded amount, and that the surgery time is controlled by the initial data.

\textbf{Finite extinction time.} Perelman proved that for any closed 3-manifold not containing an incompressible torus, the Ricci flow with surgery becomes extinct in finite time. For a simply connected 3-manifold, this implies that the only possible limit geometries are 3-spheres, establishing the Poincar\'e conjecture. The finite extinction time argument uses the monotone quantity $\mathcal{L}$-length, defined as the infimum of $\int_0^T \sqrt{\tau(t)} \int_M (|\dot{g}|^2 + R^2) \, dV \, dt$ over paths of metrics, and shows that the $\mathcal{L}$-length grows at most linearly in the number of surgeries.

The reduced volume $\tilde{V}(\bar{\tau}) = \int_M \bar{\tau}^{-n/2} e^{-\mathcal{L}(q,\bar{\tau})} \, dV$, where $\mathcal{L}$ is the $\mathcal{L}$-length functional, provides an independent monotone quantity essential for the canonical neighborhood theorem and the analysis of high-curvature regions. Perelman's key insight was that $\tilde{V}$ is monotonically increasing along the Ricci flow with surgery, providing a uniform lower bound on the reduced volume at each surgery time.

\section{Impact of This Discovery}

Perelman's proof resolved two of the most important problems in mathematics: the Poincar\'e conjecture (1904--2003) and the Thurston geometrization conjecture (1982--2003). With geometrization proved, the classification of closed 3-manifolds is now complete---every closed 3-manifold can be decomposed into geometric pieces, each carrying one of the eight Thurston geometries.

The techniques Perelman developed---the entropy functional, the no-local-collapsing theorem, the surgery construction, and the canonical neighborhood theorem---have become fundamental tools in geometric analysis. They have been applied to the differentiable sphere theorem, the classification of 4-manifolds, and the K\"ahler-Ricci flow. The Perelman entropy functional $\mathcal{W}(g,f,\tau)$ and its monotonicity formula have been generalized to other geometric evolution equations, including the Ricci flow on K\"ahler manifolds and the mean curvature flow of submanifolds in ambient spaces with curvature bounds.

The proof also established new paradigms in geometric PDEs: monotone quantities as tools for controlling geometric evolution equations, blow-up analysis using singularity models, and surgery constructions for continuing flows past singularities. These ideas have influenced the study of mean curvature flow, the harmonic map heat flow, and the Yang--Mills heat flow.

The \emph{Poincar\'e conjecture} was originally formulated for smooth manifolds: a closed, simply connected 3-manifold is homeomorphic to $S^3$. Perelman's proof actually establishes the stronger result that the manifold is \emph{diffeomorphic} to $S^3$, since the Ricci flow with surgery produces a smooth structure. The smooth Poincar\'e conjecture in dimension 3 was previously known only for manifolds with additional structure (e.g., the smooth Poincar\'e conjecture for simply connected spin 3-manifolds follows from the work of Freedman and others).

The \emph{Thurston geometrization conjecture} decomposes a closed orientable 3-manifold $M$ along incompressible tori into geometric pieces. Perelman's proof proceeds by running the Ricci flow with surgery: the surgery times correspond to the decomposing tori, and the limit geometries of the flow pieces are the eight Thurston geometries. The \emph{JSJ decomposition} (Jaco--Shalen--Johannson) provides the topological framework for the decomposition, and the Ricci flow with surgery provides the analytical framework for proving that the decomposition is canonical.

The \emph{canonical neighborhood theorem} describes the geometry of a ball $B(p, r)$ centered at a point $p$ with $R(p) \ge k^2$ and $r \le \varepsilon k^{-1}$: the ball is diffeomorphic to either (a) a round sphere, (b) a neck $S^2 \times I$, or (c) a cap (a neighborhood of a point on $S^3$ or $\mathbb{RP}^3$). The theorem is proved using the no-local-collapsing theorem and the classification of singularity models. The canonical neighborhood theorem is the key input for the surgery construction, as it ensures that all high-curvature regions have a simple geometric structure that can be controlled.

The \emph{singularity models} for the Ricci flow on 3-manifolds are solutions to the ancient Ricci flow that arise as blow-down limits of the original flow. Perelman showed that the only singularity models are: (a) the shrinking round sphere $S^3$, (b) the shrinking round cylinder $S^2 \times \mathbb{R}$, and (c) the shrinking round cylinder with a cap (the ``capped neck''). This classification is crucial for the surgery construction, as it ensures that the high-curvature regions are always necks that can be cut and capped.

The verification of Perelman's proof was carried out by several teams (Kleiner--Lott, Morgan--Tian, Cao--Zhu), producing hundreds of pages of expository notes that became standard references in geometric analysis. The total verification effort exceeded 1000 pages, and the detailed accounts provide extensive exposition of Perelman's techniques. Each verification team had to fill in details that Perelman left implicit, and the process of verification revealed the depth and subtlety of his arguments. The Kleiner--Lott notes, posted on the arXiv in 2003--2006, were the first complete account and remain the most widely cited. Morgan and Tian's monograph, published by the Clay Mathematics Institute in 2006, provided a streamlined version that emphasized the geometric analysis. Cao and Zhu's account, published by the AMS in 2006, gave a self-contained presentation suitable for graduate students.

The \emph{reduced volume} $\tilde{V}(\bar{\tau})$ is defined as the integral $\tilde{V}(\bar{\tau}) = \int_M \bar{\tau}^{-n/2} e^{-\mathcal{L}(q,\bar{\tau})} \, dV$, where $\mathcal{L}$ is the $\mathcal{L}$-length functional measuring the weighted length of paths in the space of metrics. The reduced volume satisfies $\tilde{V}(\bar{\tau}) \le (4\pi \bar{\tau})^{-n/2} \operatorname{Vol}(M)$, with equality at the initial time, and the monotonicity $\frac{d}{d\bar{\tau}} \tilde{V} \le 0$ follows from the Cauchy--Schwarz inequality applied to the gradient of $\mathcal{L}$. The reduced volume provides a uniform lower bound on the geometry of high-curvature regions, and it is the key tool for the canonical neighborhood theorem. The monotonicity of $\tilde{V}$ is equivalent to the \emph{Perelman entropy inequality}, which states that the normalized entropy $W(g,f,\tau) = -\log \tilde{V}(\bar{\tau}) + n \log(4\pi)$ is non-decreasing along the Ricci flow.

The \emph{canonical neighborhood theorem} describes the geometry of a ball $B(p, r)$ centered at a point $p$ with $R(p) \ge k^2$ and $r \le \varepsilon k^{-1}$: the ball is diffeomorphic to either (a) a round sphere, (b) a neck $S^2 \times I$, or (c) a cap (a neighborhood of a point on $S^3$ or $\mathbb{RP}^3$). The theorem is proved using the no-local-collapsing theorem and the classification of singularity models. The key step is the \emph{$\kappa$-noncollapsed condition}: a solution of the Ricci flow is $\kappa$-noncollapsed at scale $r$ if whenever $R(p) \le r^{-2}$, the volume of $B(p,r)$ is at least $\kappa r^n$. Perelman shows that if the flow is $\kappa$-noncollapsed at some scale, then the canonical neighborhood theorem holds with constants depending only on $\kappa$.

The \emph{singularity models} for the Ricci flow on 3-manifolds are ancient solutions to the Ricci flow that arise as blow-down limits of the original flow. An ancient solution is a solution defined for all $t \in (-\infty, T)$ that develops a singularity at time $T$. Perelman showed that the only singularity models are: (a) the shrinking round sphere $S^3$, (b) the shrinking round cylinder $S^2 \times \mathbb{R}$, and (c) the shrinking round cylinder with a cap (the ``capped neck''). This classification is crucial for the surgery construction, as it ensures that the high-curvature regions are always necks that can be cut and capped. The classification is proved using the reduced volume: each singularity model has a specific reduced volume asymptotics, and the monotonicity of $\tilde{V}$ forces the singularity model to be one of the three types.

\section{Modern Views and Improvements}

The Ricci flow has been applied to K\"ahler geometry, particularly in the study of K\"ahler-Einstein metrics. The \emph{K\"ahler-Ricci flow} has been used to prove the existence of K\"ahler-Einstein metrics on Fano manifolds with vanishing Futaki invariant (Tian--Donaldson--Chen--Donaldson), building directly on Perelman's entropy methods. The convergence of the K\"ahler-Ricci flow to a K\"ahler-Einstein metric was established by Tian and Chen--Donaldson independently, using Perelman's functional.

The surgery construction provides a model for handling singularities in other geometric flows, including the mean curvature flow and the Ricci flow on higher-dimensional manifolds. The classification of singularity models---neck pinches, spherical regions---has become a standard technique in the study of geometric evolution equations. Recent work by Brendle and Neves on the differentiable sphere theorem uses Perelman's canonical neighborhood theorem as a starting point.

The \emph{K\"ahler-Ricci flow} on Fano manifolds has been studied extensively using Perelman's methods. Tian proved that if the Futaki invariant vanishes and the Mabuchi energy is proper, then the K\"ahler-Ricci flow converges to a K\"ahler-Einstein metric. Chen and Donaldson gave independent proofs using Perelman's functional, and the result establishes a far-reaching generalization of the Yau--Tian--Donaldson conjecture for K\"ahler-Einstein metrics on Fano manifolds.

The \emph{mean curvature flow} of surfaces in $\mathbb{R}^3$ develops singularities that are modeled on shrinking spheres and cylinders (neck pinches). The surgery construction for mean curvature flow, developed by Brendle and Huisken, follows Perelman's pattern: identify the high-curvature region, cut it out, cap the boundaries, and restart the flow. This construction has been used to classify surfaces of mean curvature flow and to prove the differentiable sphere theorem.

The \emph{differentiable sphere theorem} (Brendle--Schoen, 2009) states that a Riemannian manifold with pointwise $1/4$-pinched sectional curvature is diffeomorphic to a space form. The proof uses the Ricci flow with surgery and Perelman's canonical neighborhood theorem, combined with a delicate analysis of the Weyl curvature tensor under the flow. The theorem was previously known only in the topological category (due to Berger), and the differentiable version requires Perelman's techniques.

The \emph{Hamilton--Perelman monotonicity formula} for the Ricci flow states that the normalized entropy $W(g,f,\tau)$ is non-decreasing along the coupled Ricci flow and backward heat equation. This monotonicity formula generalizes the classical monotonicity formulas for the mean curvature flow (Ilmanen) and the harmonic map heat flow (Struwe). The monotonicity formula is the key tool for the classification of singularity models and the finite extinction time argument.

The \emph{strong maximum principle} for the Ricci flow, proved by Hamilton, states that if the curvature tensor is non-negative at an initial time and non-zero at some point, then the curvature is strictly positive for all positive times. This principle is used in the proof of the canonical neighborhood theorem: if a region has non-negative curvature at some time, then the strong maximum principle ensures that the curvature remains non-negative, which is essential for the classification of singularity models.

The \emph{Hamilton--Ivey pinch estimate} provides a lower bound on the scalar curvature in terms of the maximum curvature: $R \ge -C |{\rm Rm}|$ for some constant $C$ depending on the initial data. This estimate prevents the curvature from becoming too negative during the Ricci flow, and it is used in the classification of singularity models. The Hamilton--Ivey estimate is proved using the evolution equation for the scalar curvature under the Ricci flow: $\partial_t R = \Delta R + 2|{\rm Rm}|^2$, combined with a maximum principle argument.

The \emph{Perelman--Burago--Gromov theorem} on Alexandrov spaces states that the tangent cone of an Alexandrov space at any point is unique, and that the tangent cone varies continuously with the base point. This result is essential for the compactness theory of Alexandrov spaces and for the surgery construction in Perelman's proof. The uniqueness of tangent cones is proved using the Gromov--Hausdorff convergence of rescalings of the space, and the continuity is proved using the monotonicity of the reduced volume.

The \emph{$\kappa$-noncollapsed condition} is a key technical tool in Perelman's proof: a solution of the Ricci flow is $\kappa$-noncollapsed at scale $r$ if whenever $R(p) \le r^{-2}$, the volume of $B(p,r)$ is at least $\kappa r^n$. Perelman shows that if the flow is $\kappa$-noncollapsed at some scale, then the canonical neighborhood theorem holds with constants depending only on $\kappa$. The $\kappa$-noncollapsed condition is preserved by the Ricci flow with surgery, and it provides the uniform control on the geometry of high-curvature regions that is needed for the surgery construction.

Perelman's earlier work on Alexandrov spaces (with Burago and Gromov) established the foundations of metric geometry, and his comparison geometry results continue to be standard references. The Grigorchuk--Perelman group, constructed with Rostislav Grigorchuk, provided examples of groups with intermediate growth in geometric group theory.

The \emph{Grigorchuk group} $G$ is the first example of a group with intermediate growth---growing faster than any polynomial but slower than any exponential function. Specifically, the growth function $a_n$ of $G$ satisfies $C_1 n^{\alpha} \le a_n \le C_2 n^{\beta}$ for $0 < \alpha < \beta < \infty$, where the exact growth rate is $n^{(\log 2)/(\log(16/3))} \approx n^{0.767}$. The construction uses a recursive wreath product construction: $G = \langle a, b, c, d \rangle$ where $a, b, c, d$ are automorphisms of the binary rooted tree satisfying specific commutation relations. The growth of $G$ is governed by the recurrence $a_n = a_{n-1}^2 + a_{n-2}^2$, which interpolates between polynomial and exponential growth. The Grigorchuk group also has the property of being a \emph{just-infinite} group: every non-trivial normal subgroup has finite index. The construction of the Grigorchuk group, together with Perelman's work on the geometrization conjecture, established that the growth of finitely generated groups is a rich invariant that interacts with both geometry and dynamics.

Perelman's work on \emph{Alexandrov spaces} (with Burago and Gromov) established the foundations of metric spaces with curvature bounded below in the sense of Alexandrov. These spaces generalize Riemannian manifolds with lower curvature bounds, and they arise naturally as Gromov--Hausdorff limits of sequences of manifolds with uniform lower curvature bounds. The theory of Alexandrov spaces provides the framework for understanding the limits of the Ricci flow with surgery, and Perelman's comparison geometry results continue to be standard references in the field. The key result is the \emph{Perelman--Burago--Gromov theorem}, which shows that Alexandrov spaces have a well-defined notion of tangent cones at every point, and that these tangent cones vary continuously. This result is essential for the compactness theory of Alexandrov spaces and for the surgery construction in Perelman's proof of the geometrization conjecture.

Perelman declined the Fields Medal in 2006 and the Clay Millennium Prize of \$1 million in 2010, withdrawing from mathematics entirely. He resigned from the Steklov Institute in 2005 and has lived in seclusion in St.~Petersburg ever since. In declining the Millennium Prize, Perelman stated that Hamilton's contribution was equal to his own, and that the prize was unfairly awarded to him alone. His refusal of both prizes has become one of the most discussed events in modern mathematics, raising questions about the nature of mathematical credit and the relationship between individual achievement and collective effort.

The \emph{$\mathcal{L}$-length functional} is defined as the infimum of $\int_0^T \sqrt{\tau(t)} \int_M (|\dot{g}|^2 + R^2) \, dV \, dt$ over paths of metrics $(g(t))_{t \in [0,T]}$ with $g(0) = g_0$ and $\tau(t)$ a time-parameterization. Perelman showed that the $\mathcal{L}$-length satisfies a monotonicity formula along the Ricci flow: if $g(t)$ is a solution of the Ricci flow, then the $\mathcal{L}$-length of any path decreases monotonically. This monotonicity is the key to the finite extinction time argument: each surgery reduces the $\mathcal{L}$-length by a bounded amount, and the total number of surgeries is finite because the $\mathcal{L}$-length is bounded below. The finite extinction time then follows from the fact that the Ricci flow on a simply connected 3-manifold must become extinct in finite time, since the only possible limit geometry is the 3-sphere.

The \emph{DeTurck trick} is a modification of the Ricci flow that makes the equation strictly parabolic. The standard Ricci flow $\partial_t g_{ij} = -2R_{ij}$ is not strictly parabolic because of diffeomorphism invariance: if $\phi_t$ is a family of diffeomorphisms, then $\phi_t^* g(t)$ is also a solution. The DeTurck modification replaces the Ricci flow with $\partial_t g_{ij} = -2R_{ij} + \nabla_i V_j + \nabla_j V_i$, where $V^k = g^{pq}(\Gamma^k_{pq} - \tilde{\Gamma}^k_{pq})$ is the DeTurck vector field and $\tilde{\Gamma}$ is a fixed background connection. The modified equation is strictly parabolic and has a unique short-time solution. The original Ricci flow is then recovered by composing with the flow of the DeTurck vector field. Perelman used the DeTurck trick to establish short-time existence and uniqueness of the Ricci flow with surgery.

\begin{center}\rule{0.5\linewidth}{0.5pt}\end{center}

\section*{Summary}
\addcontentsline{toc}{section}{Summary}

Perelman proved the Poincar\'e conjecture and the Thurston geometrization conjecture by completing Hamilton's Ricci flow program through the introduction of the entropy functional, the no-local-collapsing theorem, a surgery construction, and a finite extinction time argument. He declined both the Fields Medal and the Clay Millennium Prize.
